% ==============================================================================
% Standalone attachment: Python analysis pipeline
% ==============================================================================
\documentclass[
  a4paper, 10pt,
  master,
  final,
  notitlepage,
  noheader,
  nofrontmatter,
  noacro,
  nobibliography,
  nosoftware,
  nobackmatter,
  CTUend
]{CTUthesis}

\title{Dokumentace analytické pipeline v Pythonu}
\faculty{Masarykův ústav vyšších studií}
\department{Projektové řízení inovací}

\usepackage[most]{tcolorbox}
\tcbuselibrary{breakable,listings,skins}
\usepackage{longtable}

% Reuse the paper/colour logic already used by the poster class: white paper,
% CTU blue for structure, neutral grey for secondary text, light grey-blue fills.
\definecolor{CTUblue}{RGB}{0,101,189}
\definecolor{CTUgrey}{RGB}{84,88,90}
\definecolor{CTUlight}{RGB}{235,238,243}
\pagecolor{white}
\color{black}

\renewcommand\baselinestretch{1.08}
\setlength\parindent{0pt}
\setlength\parskip{0.45\baselineskip}
\setlist[itemize]{leftmargin=*,nosep}

\newtcblisting{pipelinecode}[1][]{%
  enhanced,
  breakable,
  listing only,
  listing engine=listings,
  colback=CTUlight!55,
  colframe=CTUblue!75!black,
  boxrule=0.45pt,
  arc=2pt,
  left=2mm,right=2mm,top=1mm,bottom=1mm,
  title={#1},
  colbacktitle=CTUblue,
  coltitle=white,
  fonttitle=\sffamily\bfseries\footnotesize,
  listing options={%
    basicstyle=\ttfamily\scriptsize,
    columns=fullflexible,
    keepspaces=true,
    breaklines=true,
    showstringspaces=false,
    frame=none,
    tabsize=2,
    upquote=true
  }
}

\tikzset{
  pipebox/.style={
    draw=CTUblue!70!black,
    fill=CTUlight,
    rounded corners=2pt,
    align=center,
    inner sep=4pt,
    minimum height=9mm,
    text width=29mm,
    font=\sffamily\scriptsize
  },
  storebox/.style={
    pipebox,
    draw=CTUgrey,
    fill=white
  },
  flow/.style={-{Latex[length=2.2mm]},line width=0.45pt,draw=CTUblue!80!black},
  weakflow/.style={-{Latex[length=2mm]},line width=0.35pt,draw=CTUgrey!80,dashed}
}

\begin{document}
\pagestyle{plain}
\small

{\sffamily\bfseries\LARGE\color{CTUblue} Dokumentace analytické pipeline v Pythonu\par}
{\sffamily\color{CTUgrey} Samostatná příloha k diplomové práci \textit{Sociální dialog a kolektivní vyjednávání}.\par}
\vspace{0.4\baselineskip}
\hrule height 0.5pt
\vspace{0.8\baselineskip}

\section*{Účel a hranice roury \textit{(pipeline)}}
Python část repozitáře vytváří reprodukovatelné obrázky, tabulky a krátké
\LaTeX{}ové vstupy používané v hlavním textu práce. Roura stahuje data z~otevřených databází
a~připravuje výstupy pro řetězec sestavení \textit{(build)} dokumentu práce. Spojuje zdrojová
data, lokální mezipaměť \textit{(cache)}, jednotlivé analýzy, registr výstupů
a~\LaTeX{}ový dokument. Ruční zásahy patří do verzovaných komentářů a~obalů
\textit{(wrappers)} PGF obrázků, nikoli do generovaných souborů
v~\texttt{latex/texparts/python/}.

\begin{center}
\resizebox{\linewidth}{!}{%
\begin{tikzpicture}[node distance=6mm and 7mm]
  \node[pipebox] (sources) {Eurostat\\OECD\\MPSV};
  \node[pipebox,right=of sources] (fetch) {\texttt{stattool.fetch}\\download + cache};
  \node[storebox,right=of fetch] (data) {\texttt{python/data/}\\raw cache};
  \node[pipebox,right=of data] (analysis) {\texttt{analyses/*.py}\\clean + compute};
  
  \node[pipebox,below=of analysis] (style) {\texttt{stattool.style}\\PDF/PGF export};
  \node[storebox,left=of style] (outputs) {\texttt{python/figures/}\\\texttt{texparts/python/}};
  \node[pipebox,left=of outputs] (quality) {\texttt{stats\_analytics.py}\\data-quality audit};
  \node[storebox,left=of quality] (reports) {\texttt{review/\\data\_quality\_report.*}};
  \node[pipebox,below=of quality] (latex) {\texttt{main.tex}\\commentary};

  \draw[flow] (sources) -- (fetch);
  \draw[flow] (fetch) -- (data);
  \draw[flow] (data) -- (analysis);
  \draw[flow] (analysis) -- (style);
  \draw[flow] (style) -- (outputs);
  \draw[flow] (outputs) -- (quality);
  \draw[flow] (quality) -- (reports);
  \draw[weakflow] (quality) -- (latex);
\end{tikzpicture}}
\end{center}

Každý analytický skript stahuje data přes projektové pomocné funkce (\textit{helpery}),
výstupy zapisuje do cest z~\texttt{config.py} a~před
vykreslováním aplikuje jednotný styl. Tím je zajištěno, že změna šablony,
barevné palety nebo způsobu exportu nevyžaduje úpravy v~jednotlivých souborech nebo hlavním textu.

\section*{Řízení generování}
Orchestrátor \texttt{python/stats\_analytics.py} nejprve rekurzivně prochází
\texttt{latex/main.tex} a~vyhledává v~něm aktivní \texttt{\textbackslash input} souborů.
Zachytává kořeny \textit{(stems)} vložené jako \texttt{texparts/python/STEM} nebo
\texttt{\textbackslash inputpgffigure\{STEM\}}. Potom porovná nalezené stemmy 
s~registrem \texttt{python/analytics\_registry.toml}; spustí pouze ty skupiny,
které jsou v~textu skutečně použité, a~kterým chybí očekávaný výstup, případně
ty, které byly vynuceny volbou \texttt{--force}. Po dokončení generování navíc
provede audit datové kvality -- vypíše varování pro záložní (\textit{fallback}) cesty,
zkontroluje uvedený rok v~popiscích a~datech, a~upozorní na použití jiného roku než je
cílový.

\begin{center}
\resizebox{0.92\linewidth}{!}{%
\begin{tikzpicture}[node distance=7mm and 10mm]
  \node[pipebox,text width=34mm] (scan) {scan \texttt{main.tex}\\a vložených texpartů};
  \node[pipebox,text width=34mm,right=of scan] (stems) {aktivní stemmy\\\texttt{texparts/python}};
  \node[pipebox,text width=34mm,right=of stems] (registry) {match v\\\texttt{analytics\_registry.toml}};
  \node[pipebox,text width=34mm,right=of registry] (run) {spustit skript\\jen když chybí výstup};
  \node[pipebox,text width=34mm,below=of run] (audit) {audit roku dat\\a fallback varování};
  \node[storebox,text width=34mm,below=of registry] (skip) {neaktivní nebo kompletní\\skupiny se přeskočí};
  \node[storebox,text width=34mm,left=of audit] (report) {\texttt{review/data\_quality\_report.*}};

  \draw[flow] (scan) -- (stems);
  \draw[flow] (stems) -- (registry);
  \draw[flow] (registry) -- (run);
  \draw[flow] (run) -- (audit);
  \draw[flow] (audit) -- (report);
  \draw[weakflow] (registry) -- (skip);
\end{tikzpicture}}
\end{center}

\begin{pipelinecode}[Běžné spuštění]
cd python
bash run.sh stats_analytics.py              # doplni chybejici aktivni vystupy
bash run.sh stats_analytics.py --list       # vypise stemmy nalezene v LaTeXu
bash run.sh stats_analytics.py --force all  # pregeneruje vsechny aktivni skupiny
bash run.sh stats_analytics.py --target-year 2025  # explicitni audit proti cili 2025
\end{pipelinecode}

Datová kvalita je od května 2026 součástí stejného kroku, nikoli samostatná
ruční kontrola. Výchozí cílový rok je 2025; lze jej změnit volbou
\texttt{--target-year} nebo proměnnou prostředí \texttt{DP\_TARGET\_YEAR}.
Analytické skripty mají používat helpery \texttt{warn\_fallback(...)}
a~\texttt{warn\_non\_target\_year(...)} z~modulu \texttt{stattool.data\_quality},
takže se do logu promítne každý sekundární zdroj, hardcoded hodnota, expertní
odhad i~použití jiného roku než 2025.

Výstupem jsou varování v~terminálu, \LaTeX{} tabulka (vložená níže) a~dva soubory
pro transparentní deklaraci v~popiscích a komentářích:

\begin{itemize}
  \item \texttt{review/data\_quality\_report.json}: strojově čitelný přehled auditovaných stemmů.
  \item \texttt{review/data\_quality\_report.md}: stručný report pro lidskou revizi.
\end{itemize}

\section*{Smlouva mezi Pythonem a LaTeXem}
Registr je jediným místem, kde se analytický skript mapuje na výstupy. Klíč
seskupuje logicky související grafy nebo tabulky; atribut \texttt{texparts}
určuje stemmy hledané v~\LaTeX{}u a~atribut \texttt{pics} slouží pro kontrolu
PDF výstupu. PGF výstupy mají navíc v~git verzovaný \LaTeX{} soubor
\texttt{latex/texparts/figures/STEM.tex}, v~němž lze upravovat řetězce uvnitř grafu
a~posouvat popisky bez přepisu při dalším běhu Pythonu.

\begin{pipelinecode}[Příklad z registru]
[income_pps]
script   = "analyses/eu_prijem_pps.py"
texparts = ["eu_prijem_pps"]
pics     = ["eu_prijem_pps"]
\end{pipelinecode}

Každá analýza má stejnou kostru bez ohledu na zdroj dat: import cest
a~stylu, stažení dat přes pomocnou funkci, jejich zpracování balíčkem pandas,
vykreslení PDF nebo PGF a~zápis obalu v~\LaTeX{}u. V~PGF podložce
\textit{(backend)} se text v~grafu vykresluje až v~hlavním dokumentu, takže
v~popiscích fungují \LaTeX{}makra jako jsou jednotky \texttt{\textbackslash SI\{\}\{\}}
nebo zkratky \texttt{\textbackslash acs\{geo-CZ\}}.

\begin{pipelinecode}[Kostra analytického skriptu]
from config import LATEX_TEXPARTS_DIR
from stattool.fetch import fetch_eurostat
from stattool.style import apply_style_pgf, savefig_pgf, save_figure_tex_pgf

COUNTRIES = ["CZ", "AT", "DE", "DK", "PL", "SK"]

apply_style_pgf()
data = fetch_eurostat("dataset_code")
# ... filtrovani, cisteni, vypocet ukazatele ...

savefig_pgf(fig, "my_stem")
save_figure_tex_pgf(
    "my_stem",
    caption="Co, uzemni rozsah, rok",
    label="fig:my_stem",
    cite_keys=["eurostat_dataset_code"],
    strings={"title": "Titulek renderovany LaTeXem"},
)
\end{pipelinecode}

\section*{Výstupy a ruční úpravy}
Podložka PDF vytváří hotový graf a~\LaTeX{} obal,
který PDF vkládá přes \texttt{\textbackslash includegraphics}. Podložka
PGF vytváří zdrojový \texttt{.pgf} grafický soubor
v~\texttt{python/figures/}, jednořádkový obal
v~\texttt{latex/texparts/python/} a~editovatelný soubor s~přepisovatelnými řetězci
v~\texttt{latex/texparts/figures/}. Binární doprovodné obrázky PGF exportu se
deduplikují do sdílených bitmap v~\texttt{python/figures/\_shared/}. Výsledkem je, že velké
mapy mohou zůstat datově generované, ale jejich titulky a~textové řetězce lze
ladit v~\LaTeX{}u a~velikost na disku je minimální.

\begin{itemize}
  \item \texttt{python/data/}: lokální cache stažených dat, mimo git.
  \item \texttt{python/figures/}: generované PDF/PGF výstupy, mimo git.
  \item \texttt{latex/texparts/python/}: generované vstupy, mimo git.
  \item \texttt{latex/texparts/figures/}: ručně doladitelné PGF wrappery, v~gitu.
  \item \texttt{latex/texparts/commentary/}: metodologická a~interpretační próza.
\end{itemize}

\section*{Sestavení \textit{(build)} a kontrola}
Ve VS Code je běžný postup zprostředkován recepty LaTeX Workshopu: recepty
s~analýzami nejprve vytvoří python prostředí \texttt{setup\_venv}, poté spustí orchestrátor
\texttt{stats\_analytics.py} a~poté sestavení dokumentu v~\LaTeX{}u. Při ručním sestavení
lze rouru spustit explicitně nebo nastavit \texttt{RUN\_PYTHON\_ANALYTICS=1} jako
háček před sestavením \textit{(pre-build hook)} přes \texttt{latexmk}.

\begin{pipelinecode}[Ruční build]
cd latex
RUN_PYTHON_ANALYTICS=1 latexmk -pdf -interaction=nonstopmode -file-line-error -outdir=build main.tex
pdflatex -synctex=1 -interaction=nonstopmode -file-line-error -output-directory=build python_pipeline_attachment.tex
\end{pipelinecode}

Je-li recept spuštěn z~VS Code, zápis (\textit{log}) z~orchestrace se současně
ukládá do \texttt{latex/build/stats-analytics.log} (případně
\texttt{stats-analytics-force.log}). Jsou v~něm zachyceny všechny varování a~chyby, včetně
záložních datových cest a~upozornění na použitý rok. Stejná upozornění jsou za běhu
vytištěna v~terminálu nebo uložena ve zprávách v~adresáři \texttt{review/}.

Chybějící \texttt{.pgf} nebo \texttt{texparts/python} výstup je třeba řešit spuštěním orchestrátoru,
chyba v~citaci nebo makru se řeší v~\texttt{texparts/figures}. Ruční úprava řetězců se za běhu orchestrátoru
zapíše do příslušného python skriptu z~editovatelného obalu PGF.

\section*{Tabulka datové kvality}
Následující přehled byl vygenerován při běhu \texttt{stats\_analytics.py} při sestavení práce.

\IfFileExists{../review/data_quality_pipeline_table.tex}{%
  % Auto-generated by python/stats_analytics.py. Do not edit manually.
\begingroup
\scriptsize
\setlength{\tabcolsep}{3pt}
\renewcommand{\arraystretch}{1.12}
\begin{longtable}{p{0.08\linewidth}p{0.20\linewidth}p{0.38\linewidth}p{0.29\linewidth}}
\hline
Obr. & Skript & Dataset(y) a filtry & Pozn\'{a}mky (roky, fallback) \\
\hline
\endfirsthead
\hline
Obr. & Skript & Dataset(y) a filtry & Pozn\'{a}mky (roky, fallback) \\
\hline
\endhead
fig:eu\_apz\_vydaje\_2004 & eu\_\allowbreak{}apz\_\allowbreak{}vydaje.py & OECD LMPEXP (via ``load\_lmp\_active()``) & Roky v popisku: 2004, 2023 | Nesedi cilovy rok 2025 \\
fig:eu\_bohatstvi\_top1\_wid & eu\_\allowbreak{}bohatstvi\_\allowbreak{}top1\_\allowbreak{}wid.py & World Inequality Database (WID) bulk CSV: https://wid.\allowbreak{}world/bulk\_download/ & Roky v popisku: 2024 | Nesedi cilovy rok 2025 \\
fig:eu\_cenova\_hladina & eu\_\allowbreak{}cenova\_\allowbreak{}hladina.py & Eurostat ``prc\_ppp\_ind``; Eurostat prc\_ppp\_ind (A.\allowbreak{}PLI\_EU27\_2020.\allowbreak{}GDP.\allowbreak{}, start\_period=START\_YEAR) & Roky v popisku: 2024 | Nesedi cilovy rok 2025 \\
fig:eu\_danovy\_klin & eu\_\allowbreak{}danovy\_\allowbreak{}klin.py & Eurostat, earn\_nt\_taxwedge (start\_period=2015); Eurostat earn\_nt\_taxwedge (start\_period=2015) & Roky v popisku: 2025 \\
fig:eu\_gini\_prijem & eu\_\allowbreak{}gini\_\allowbreak{}prijem.py & Eurostat, ``ilc\_di12``; Eurostat ilc\_di12 (A.\allowbreak{}TOTAL.\allowbreak{}GINI\_HND.\allowbreak{}, start\_period=START\_YEAR) & Roky v popisku: 2003, 2025 \\
fig:eu\_hustota\_mapa & eu\_\allowbreak{}hustota\_\allowbreak{}mapa.py & OECD AIAS ICTWSS (dataset ``TUD``) & Roky v popisku: 2024 | Nesedi cilovy rok 2025 \\
fig:eu\_ict\_gender & eu\_\allowbreak{}ict\_\allowbreak{}gender.py & Eurostat ``isoc\_sks\_itsps``; Eurostat isoc\_sks\_itsps (start\_period=2015) & Roky v popisku: 2025 \\
fig:eu\_mne\_mapa & eu\_\allowbreak{}mne\_\allowbreak{}mapa.py & Eurostat ``egr\_emp``; Eurostat egr\_emp (start\_period=2015) & Roky v popisku: 2024 | Nesedi cilovy rok 2025 \\
fig:eu\_odvetvove\_mzdy\_bar & eu\_\allowbreak{}odvetvove\_\allowbreak{}mzdy.py & Eurostat lc\_lci\_lev (f"A.\allowbreak{}EUR.\allowbreak{}D1\_D4\_MD5.\allowbreak{}\{sector\_filter\}.\allowbreak{}", start\_period=DISPLAY\_YEAR); Eurostat prc\_ppp\_ind (A.\allowbreak{}PLI\_EU27\_2020.\allowbreak{}GDP.\allowbreak{}, start\_period=DISPLAY\_YEAR) & Roky v popisku: 2025 | Fallback vetve: f"EU27 sector benchmark missing for \{', '.join(\_fallback\_sectors)\}; cross-country mean used" \\
fig:eu\_odvetvove\_mzdy\_mapa & eu\_\allowbreak{}odvetvove\_\allowbreak{}mzdy.py & Eurostat lc\_lci\_lev (f"A.\allowbreak{}EUR.\allowbreak{}D1\_D4\_MD5.\allowbreak{}\{sector\_filter\}.\allowbreak{}", start\_period=DISPLAY\_YEAR); Eurostat prc\_ppp\_ind (A.\allowbreak{}PLI\_EU27\_2020.\allowbreak{}GDP.\allowbreak{}, start\_period=DISPLAY\_YEAR) & Roky v popisku: 2025 \\
fig:eu\_osvc\_mapa & eu\_\allowbreak{}osvc\_\allowbreak{}mapa.py & Eurostat lfsa\_egaps (Employed persons by professional status); Eurostat lfsa\_egaps (A.\allowbreak{}THS\_PER.\allowbreak{}T.\allowbreak{}Y15-74.\allowbreak{}EMP+SELF.\allowbreak{}, start\_period=2015) & Roky v popisku: 2025 \\
fig:eu\_pokryti\_kv\_mapa & eu\_\allowbreak{}pokryti\_\allowbreak{}kv\_\allowbreak{}mapa.py & ICTWSS AdjCov (adjusted coverage) for most EU-27 countries; & Roky v popisku: 2024 | Nesedi cilovy rok 2025 \\
fig:eu\_prijem\_distribuce & eu\_\allowbreak{}prijem\_\allowbreak{}distribuce.py & Eurostat ilc\_di01 (f"A.\allowbreak{}\{quant\_filter\}.\allowbreak{}TC.\allowbreak{}PPS.\allowbreak{}\{geo\_filter\}", start\_period=START\_YEAR) & Roky v popisku: 2025 \\
fig:eu\_produktivita\_prijem\_trajektorie & eu\_\allowbreak{}produktivita\_\allowbreak{}prijem\_\allowbreak{}trajektorie.py & Eurostat nama\_10\_lp\_ulc (A.\allowbreak{}PC\_EU27\_2020\_MPPS\_CP.\allowbreak{}NLPR\_HW.\allowbreak{}, start\_period=START\_YEAR); Eurostat earn\_nt\_net (A.\allowbreak{}PPS.\allowbreak{}NET.\allowbreak{}P1\_NCH\_AW100.\allowbreak{}, start\_period=START\_YEAR); Eurostat lfsa\_ewhan2 (A.\allowbreak{}TOTAL.\allowbreak{}EMP.\allowbreak{}TOTAL.\allowbreak{}Y15-64.\allowbreak{}T.\allowbreak{}HR.\allowbreak{}, start\_period=START\_YEAR) & Roky v popisku: 2008, 2025 \\
fig:korelace\_scatter & korelace\_\allowbreak{}analyza.py & Eurostat lfsa\_ewhan2 (A.\allowbreak{}TOTAL.\allowbreak{}EMP.\allowbreak{}TOTAL.\allowbreak{}Y15-64.\allowbreak{}T.\allowbreak{}HR.\allowbreak{}, start\_period=START\_YEAR); Eurostat earn\_nt\_net (A.\allowbreak{}PPS.\allowbreak{}NET.\allowbreak{}P1\_NCH\_AW100.\allowbreak{}, start\_period=START\_YEAR); Eurostat earn\_gr\_gpgr2 (A.\allowbreak{}PC.\allowbreak{}B-S\_X\_O.\allowbreak{}, start\_period=START\_YEAR); Eurostat ilc\_di12 (A.\allowbreak{}TOTAL.\allowbreak{}GINI\_HND.\allowbreak{}, start\_period=START\_YEAR); Eurostat nama\_10\_lp\_ulc (A.\allowbreak{}PC\_EU27\_2020\_MPPS\_CP.\allowbreak{}NLPR\_HW.\allowbreak{}, start\_period=START\_YEAR) & Roky v popisku: 2024 | Nesedi cilovy rok 2025 \\
fig:practical\_ternary\_social\_dialog & practical\_\allowbreak{}ternary\_\allowbreak{}social\_\allowbreak{}dialog.py & Neuvedeno (sdilene helpery nebo staticka data) & Roky v popisku: 2025, 2026 | Nesedi cilovy rok 2025 \\
fig:practical\_ternary\_strength\_map & practical\_\allowbreak{}ternary\_\allowbreak{}social\_\allowbreak{}dialog.py & Neuvedeno (sdilene helpery nebo staticka data) & Rok v popisku: neuvedeno | Rok v popisku nenalezen \\
fig:problemy\_danovy\_klin\_cz & problemy\_\allowbreak{}cz\_\allowbreak{}model.py & Model podle legislativy ČR (zákon č.\allowbreak{} 586/1992 Sb.\allowbreak{} o daních z příjmů, zákon č.\allowbreak{} 155/1995 Sb.\allowbreak{} o důchodovém pojištění, nařízení vlády č.\allowbreak{} 365/2025 Sb.\allowbreak{}) & Roky v popisku: 2026 | Nesedi cilovy rok 2025 \\
fig:problemy\_dojezdeni\_nuts2 & problemy\_\allowbreak{}dojezdeni.py & Eurostat, ``lfst\_r\_lfe2ecomm``; Eurostat lfst\_r\_lfe2ecomm (start\_period=START\_YEAR) & Roky v popisku: 2025 \\
fig:problemy\_dojezdeni\_vyvoj & problemy\_\allowbreak{}dojezdeni.py & Eurostat, ``lfst\_r\_lfe2ecomm``; Eurostat lfst\_r\_lfe2ecomm (start\_period=START\_YEAR) & Roky v popisku: 2005, 2025 \\
fig:problemy\_duchod\_solidarita & problemy\_\allowbreak{}cz\_\allowbreak{}model.py & Model podle legislativy ČR (zákon č.\allowbreak{} 586/1992 Sb.\allowbreak{} o daních z příjmů, zákon č.\allowbreak{} 155/1995 Sb.\allowbreak{} o důchodovém pojištění, nařízení vlády č.\allowbreak{} 365/2025 Sb.\allowbreak{}) & Roky v popisku: 2026 | Nesedi cilovy rok 2025 \\
fig:problemy\_gpg\_sektor & problemy\_\allowbreak{}stratifikace.py & Eurostat earn\_gr\_gpgr2 (f"A.\allowbreak{}PC.\allowbreak{}B-S\_X\_O.\allowbreak{}", start\_period=START\_YEAR) & Roky v popisku: 2010, 2024 | Nesedi cilovy rok 2025 \\
fig:problemy\_gpg\_stratifikace & problemy\_\allowbreak{}gpg.py & Eurostat earn\_gr\_gpgr2 (start\_period=2018); Eurostat earn\_ses\_hourly (ses\_filter); Eurostat prc\_ppp\_ind (f"A.\allowbreak{}PLI\_EU27\_2020.\allowbreak{}GDP.\allowbreak{}\{GEO\_6\}", start\_period=2014) & Roky v popisku: 2022 | Nesedi cilovy rok 2025 | Fallback vetve: "PLI missing for exact SES year; latest available country PLI used for conversion " + ", ".join(\_countries); "PLI unavailable for some SES rows; dropping unconvertible EUR rows (no EUR fallback) for " + ", ".join(\_countries\_missing); PLI unavailable for all SES rows; percentile profile omitted because EUR→PPS conversion is impossible \\
fig:problemy\_jazyky\_isced & problemy\_\allowbreak{}jazyky.py & Eurostat Adult Education Survey (AES), \textasciitilde{}5-year rounds (2007/2011/2016/2022).\allowbreak{}; Eurostat edat\_aes\_l21; Eurostat edat\_aes\_l22; Eurostat edat\_aes\_l23 & Roky v popisku: 2022 | Nesedi cilovy rok 2025 \\
fig:problemy\_mzda\_duchod & problemy\_\allowbreak{}mzda\_\allowbreak{}duchod.py & MPSV ISPV (Mzdová sféra), ČSSZ (Statistická ročenka důchodového pojištění) & Roky v popisku: 2024, 2025 | Fallback vetve: N\_WAGE\_PRIVATE; N\_WAGE\_PUBLIC; N\_PENSION \\
fig:problemy\_regiony\_mapa & problemy\_\allowbreak{}stratifikace.py & Eurostat earn\_gr\_gpgr2 (f"A.\allowbreak{}PC.\allowbreak{}B-S\_X\_O.\allowbreak{}", start\_period=START\_YEAR) & Roky v popisku: 2025 \\
fig:problemy\_verejny\_soukromy & problemy\_\allowbreak{}verejny\_\allowbreak{}soukromy.py & MPSV ISPV / TREXIMA (Mzdová sféra MZS, Platová sféra PLS, rok 2025) & Roky v popisku: 2025 \\
fig:stav\_arope\_dual & stav\_\allowbreak{}arope\_\allowbreak{}dual.py & Eurostat ilc\_peps01n (A.\allowbreak{}PC.\allowbreak{}TOTAL.\allowbreak{}T.\allowbreak{}); Eurostat ilc\_di01 (A.\allowbreak{}D1+D2+D3+D4+D5+D6+D7+D8+D9.\allowbreak{}TC.\allowbreak{}PPS+EUR.\allowbreak{}); Eurostat prc\_ppp\_ind (A.\allowbreak{}PLI\_EU27\_2020.\allowbreak{}GDP.\allowbreak{}EU27\_2020) & Roky v popisku: 2017, 2025 | Fallback vetve: f"EU27 D3 PPS series missing in years \{fallback\_years\}; EUR converted to PPS via PLI used as fallback" \\
fig:stav\_hdp\_vyvoj & stav\_\allowbreak{}hdp\_\allowbreak{}vyvoj.py & Eurostat, nama\_10\_pc (A.\allowbreak{}PC\_EU27\_2020\_HAB\_MPPS\_CP.\allowbreak{}B1GQ.\allowbreak{}, start\_period=START\_YEAR); Eurostat nama\_10\_pc (A.\allowbreak{}PC\_EU27\_2020\_HAB\_MPPS\_CP.\allowbreak{}B1GQ.\allowbreak{}, start\_period=START\_YEAR) & Roky v popisku: 2004, 2025 \\
fig:stav\_hustota\_vyvoj\_2004 & stav\_\allowbreak{}hustota\_\allowbreak{}vyvoj.py & OECD AIAS ICTWSS (via OECD Stats API, dataset ``TUD``) & Roky v popisku: 2004, 2024 | Nesedi cilovy rok 2025 \\
fig:stav\_ipp\_mzdy & stav\_\allowbreak{}ipp\_\allowbreak{}mzdy.py & Eurostat lc\_lci\_r2\_a (f"A.\allowbreak{}I20.\allowbreak{}B-S.\allowbreak{}D1\_D4\_MD5.\allowbreak{}\{GEO\_ALL\}", start\_period=START\_YEAR - 1) & Roky v popisku: 2007, 2025 \\
fig:stav\_ipp\_rozsah & stav\_\allowbreak{}ipp\_\allowbreak{}rozsah.py & MPSV IPP (Informační systém o pracovních podmínkách, 2007--2025) & Roky v popisku: 2007, 2025 \\
fig:stav\_kolektivni\_spory\_vyvoj & stav\_\allowbreak{}kolektivni\_\allowbreak{}spory\_\allowbreak{}vyvoj.py & ČMKOS (Zpráva o průběhu kolektivního vyjednávání v ČR, počty sporů před zprostředkovatelem a rozhodcem, 2013--2025) & Roky v popisku: 2013, 2025 \\
fig:stav\_korporatni\_dan\_mapa & stav\_\allowbreak{}korporatni\_\allowbreak{}dan.py & OECD CTS\_CIT (start\_period=START\_YEAR) & Roky v popisku: 2023 | Nesedi cilovy rok 2025 \\
fig:stav\_socialni\_mir\_skore\_mapa & stav\_\allowbreak{}socialni\_\allowbreak{}mir\_\allowbreak{}mapa.py & ETUI, ILOSTAT a expertni odhady (Social-peace benchmark, 2025) & Roky v popisku: 2025 | Fallback vetve: B4 benchmark combines ETUI/ILOSTAT with national fallback assumptions for lacunae \\
fig:stav\_stavky & stav\_\allowbreak{}stavky.py & ILOSTAT STR\_DAYS\_ECO\_RT\_A; Statistics Denmark ABST1; Eurostat lfsa\_ewhun2; Eurostat lfsi\_emp\_a (A.\allowbreak{}EMP\_LFS.\allowbreak{}T.\allowbreak{}Y20-64.\allowbreak{}THS\_PER.\allowbreak{}DK, start\_period=START\_YEAR); Eurostat lfsa\_ewhun2 (A.\allowbreak{}TOTAL.\allowbreak{}EMP.\allowbreak{}TOTAL.\allowbreak{}Y15-64.\allowbreak{}T.\allowbreak{}HR.\allowbreak{}, start\_period=START\_YEAR) & Roky v popisku: 2000, 2025 | Fallback vetve: f"ILOSTAT strike dataset used fallback value column '\{obs\_col\}' instead of 'obs\_value'" \\
fig:stav\_volna\_mista\_vyvoj & stav\_\allowbreak{}volna\_\allowbreak{}mista\_\allowbreak{}vyvoj.py & Eurostat jvs\_a\_r21 & Roky v popisku: 2010, 2025 | Fallback vetve: f"Vacancy-rate timeline used \{\_chosen\_nace\} aggregate instead of preferred B-S\_X\_O"; f"Vacancy-rate timeline used \{\_chosen\_unit\} instead of preferred AVG\_3Y"; "Vacancy-rate timeline filled missing countries by relaxing size class GE10 -> TOTAL: " + ", ".join(\_size\_overrides); "Vacancy-rate timeline filled missing countries with alternate NACE aggregates: " + ", ".join(\_country\_overrides) \\
fig:stav\_zamestnanost & stav\_\allowbreak{}zamestnanost.py & Eurostat, ``lfsi\_emp\_a``; Eurostat lfsi\_emp\_a (A.\allowbreak{}EMP\_LFS.\allowbreak{}T.\allowbreak{}Y20-64.\allowbreak{}PC\_POP.\allowbreak{}, start\_period=START\_YEAR) & Roky v popisku: 2009, 2025 \\
fig:vyhled\_digit\_dovednosti\_mapa & vyhled\_\allowbreak{}digit\_\allowbreak{}dovednosti.py & Eurostat ``isoc\_sk\_dskl\_i21``; Eurostat isoc\_sk\_dskl\_i21 (A.\allowbreak{}IND\_TOTAL.\allowbreak{}I\_DSK2\_AB.\allowbreak{}PC\_IND.\allowbreak{}, start\_period=2021) & Roky v popisku: 2025 \\
fig:vyhled\_porodnost\_mapa & vyhled\_\allowbreak{}porodnost.py & Eurostat, demo\_find (A.\allowbreak{}TOTFERRT.\allowbreak{}, start\_period=START\_YEAR); Eurostat demo\_find (A.\allowbreak{}TOTFERRT.\allowbreak{}, start\_period=START\_YEAR) & Roky v popisku: 2024 | Nesedi cilovy rok 2025 \\
fig:vyhled\_porodnost\_vyvoj & vyhled\_\allowbreak{}porodnost.py & Eurostat, demo\_find (A.\allowbreak{}TOTFERRT.\allowbreak{}, start\_period=START\_YEAR); Eurostat demo\_find (A.\allowbreak{}TOTFERRT.\allowbreak{}, start\_period=START\_YEAR) & Roky v popisku: 1990, 2024 | Nesedi cilovy rok 2025 \\
fig:vyhled\_vzdelani\_vyvoj & vyhled\_\allowbreak{}vzdelani.py & Eurostat ``edat\_lfse\_03``; Eurostat edat\_lfse\_03 (A.\allowbreak{}T.\allowbreak{}Y25-64.\allowbreak{}PC.\allowbreak{}ED5-8.\allowbreak{}, start\_period=START\_YEAR) & Roky v popisku: 2004, 2025 \\
fig:vyhled\_zavislost\_mapa & vyhled\_\allowbreak{}zavislost.py & Eurostat, ``demo\_pjanind``; Eurostat demo\_pjanind (A.\allowbreak{}OLDDEP1.\allowbreak{}, start\_period=START\_YEAR) & Roky v popisku: 2025 \\
\hline
\end{longtable}
\endgroup

}{%
  \textit{Tabulka zatím není vygenerována. Spusťte: \texttt{bash run.sh stats\_analytics.py}}.
}

\end{document}